\scanpage{113}
% archive foliation: 112

\uline{Proposition} 4.2 \deleted{\unsure{a}}

a) Stabilité de ``eff$^{t}$ adm. de poids $\subset\Phi$ et niveaux $\leq\nu$'' par sous-faisceaux constr., quotients constructibles, sommes finies, extensions.

b) $F_i$ ``eff adm de poids $\subset\Phi_i$, niveaux $\leq\nu_i$'' pour $i=1,2$ implique $F_1\otimes F_2$ ``eff$^{t}$ adm. de poids $\subset\Phi_1+\Phi_2$, niveaux $\leq\nu_1+\nu_2$''.

c) Stabilité par changement de base.

d) Si $F$ \inserted{eff$^{t}$} admissible \uline{constant tordu} ``à poids \deleted{$\unsure{j}$}\inserted{$\nu+2k$} et niveaux $\leq\nu$'' (où $k\geq0,\nu\geq0$) \deleted{\unsure{et si}} alors pour tout $n\geq\nu+k$ \inserted{(soit $n=\nu+k+h$, $h\geq0$)}, $\check F(-n)$ est eff$^{t}$ admissible constant tordu, ``à poids $2n-\nu-2k=\nu+2h$, niveaux $\leq\nu$''.

\uline{Remarque}: introduire plutôt: ``à morceaux de poids $i_\alpha$, niveaux $\nu_\alpha$'', $\alpha\in\cdots$.

\uline{Théorème} 5.1. Soit $f:X\longrightarrow S$ morphisme séparé de \deleted{t.f.} \inserted{prés finie}, et $F$ sur $X$ \deleted{eff$^{t}$ admissible} \inserted{\uuline{strictement spécial}} \deleted{de poids $j$ \unsure{et niveaux} \unsure{$\nu$}} \deleted{(i.e. déf. 2.1 \unsure{sans} \unsure{de poids} $\subset\Phi$ \unsure{niveaux})} (\unsure{i.e.} déf. 1 \unsure{p108} s'applique, avec $F$ \uline{facteur direct} d'un $R^j g_*(\mathbb Q_{lY})$, où $g:Y\longrightarrow X$ \deleted{\unsure{étant}} proj. lisse). Alors $R^i_!f_{\unsure{*|!}}(F)$ est \uline{eff$^{t}$ admissible} \deleted{\uline{\unsure{assez spécial}}} à \deleted{niveaux} poids $\leq\deleted{\unsure{2}}i+j$ \deleted{\unsure{effectif}} \inserted{et niveaux $\leq\operatorname{Inf}(i+j,2n-i+j)$}.

\deleted{\unsure{poids} $\unsure{i+j}$, \unsure{niveau} $\unsure{X}$ \unsure{pour} \unsure{tout} $\unsure{i\geq n}$,}
\deleted{Si \unsure{niveau} $n=\operatorname{Sup}_{\unsure{s\in S}}\dim X_s$ et si $i\geq n$ alors les poids}
Si $n=\operatorname{Sup}_{s\in S}\dim X_s$, \deleted{$\unsure{s\in S}$} et si $i\geq n$, les poids \deleted{\unsure{et}} et niveaux \deleted{\unsure{sont}} possibles sont
\diagram{
\begin{array}{r|c|c|c|c|}
\text{poids}&2(i-n)&2(i-n)+1&\cdots&(i+j)\\ \hline
\text{niveaux}&\leq0&\leq1&\cdots&\leq2n-i+j
\end{array}
}

\margin{\uline{Utilise Lefschetz \unsure{N.} Deligne}}
\margin{$+$ résolution des singularités \inserted{forte} pour les clôtures parfaites des corps résiduels de $S$.}
\margin{$+$ $\unsure{\frac12}$ simplicité arithmétique?}

\uline{Démonstration}. On \deleted{\unsure{en}} peut supposer $F=R^j g_*(\mathbb Q_{lY})$, $g:Y\longrightarrow X$ proj. lisse. \deleted{$\unsure{F=\mathbb Q_{lX}}$}. Utilisant le th. de changement de base propre, \deleted{\unsure{on peut}} et dévissage via suite exacte de cohomologie à supports propres, on est ramené au cas où \deleted{$\exists\,S'\longrightarrow S$} $S$ est affine, \unsure{exclt}, intègre, \inserted{$X$ affine,} $\exists\,S'\longrightarrow S$ morphisme fini plat radiciel, tel que $X' = X\times_S S'$ est \uline{lisse} sur $S'$, et se ``compactifie'' en un $\widehat X_{S'}$ \deleted{tel que $X'$} projectif et lisse, tel que $\widehat X_{S'}-X_{S'}$ soit un diviseur à croisements normaux relatifs, à composantes irréductibles \uline{lisses} sur $S'$; \inserted{enfin $F$ \deleted{\unsure{sur le même}} étant \uline{spécial} de poids $j$.} Quitte à changer de notations, on peut supposer $S=S'$, et dire $X\subset X'=\widehat X$. On termine alors \inserted{si $F=\mathbb Q_{lX}$,} comme dans les notes préliminaires (partie A, 1)). [\uline{NB} On n'utilise pas la dégénérescence de la suite spectrale qui y intervient.]

\margin{à expliciter \unsure{en faisant constamment servir le principe de} 3.3 a) et 4.2 a)}
